% Dokumententyp, benutzte Pakete und deren Einstellungen
\documentclass[	\hsmasprache,%
				\hsmaabgabe,%
				\hsmapublizieren,%
				\genderhinweis,
				\hsmaquellcode,
				\hsmasymbole,
				\hsmaglossar,
				\hsmatc]{HMA}

% Wo sind die Bilder?
\graphicspath{{bilder/}}


% Wo liegt Sourcecode?
\newcommand{\srcloc}{src/}


% Checklisten mit zwei Ebenen
\newlist{checklist}{itemize}{2}
\setlist[checklist]{label=$\square$}



% Befehl zum Erstellen eigener Makros. In diesem Fall ist es ein Makro für das Einbinden von Bildern. Das label (für \ref) ist dann der Name der Bilddatei
\newcommand{\bild}[3]{
	\begin{figure}[ht]
		\centering
		\includegraphics[width=#2]{#1}
		\caption{#3}
		\label{#1}
\end{figure}}




							


\newcommand{\snowcard}[9]{
	\begin{table}[ht!]
		\caption{\hsmasnowcardanforderung\ #1 -- #4}\label{#1}
		\renewcommand{\arraystretch}{1.2}
		\centering
		\sffamily
		\begin{footnotesize}

			\begin{tabularx}{\linewidth}{sssssl}
				\toprule
				\textbf{\hsmasnowcardno} & #1 & \textbf{\hsmasnowcardart} & #2 & \textbf{\hsmasnowcardprio} & #3 \\
				\midrule
				\multicolumn{2}{l}{\textbf{\hsmasnowcardtitel}} & \multicolumn{4}{l}{\parbox[t]{11.8cm}{#4}} \\
				\ifx&#5&%
				\else
				\multicolumn{2}{l}{\textbf{\hsmasnowcardherkunft}} & \multicolumn{4}{l}{\parbox[t]{11.8cm}{#5}} \\
				\fi
				\ifx&#6&%
				\else
				\multicolumn{2}{l}{\textbf{\hsmasnowcardkonflikt}} & \multicolumn{4}{l}{\parbox[t]{11.8cm}{#6}} \\
				\fi
				\addlinespace
				\multicolumn{6}{l}{\textbf{\hsmasnowcardbeschreibung}} \\
				\multicolumn{6}{l}{\parbox[t]{13.5cm}{#7\strut}} \\
				\ifx&#8&%
				\else
				\addlinespace
				\multicolumn{6}{l}{\textbf{\hsmasnowcardfitkriterium}} \\
				\multicolumn{6}{l}{\parbox[t]{13.5cm}{#8\strut}} \\
				\fi
				\ifx&#9&%

				\else
				\addlinespace
				\multicolumn{6}{l}{\textbf{\hsmasnowcardmaterial}} \\
				\multicolumn{6}{l}{\parbox[t]{13.5cm}{#9\strut}} \\
				\fi
				\bottomrule
			\end{tabularx}
		\end{footnotesize}
	\end{table}
}



% Quality Attribute Scenario
\newcommand{\qas}[9]{
	\begin{table}[ht!]
		\caption{\hsmaqasanforderung\ #1 -- #3}\label{#1}
		\renewcommand{\arraystretch}{1.2}
		\centering
		\sffamily
		\begin{footnotesize}

			\begin{tabularx}{\linewidth}{sssssl}
				\toprule
				\textbf{\hsmaqasno} & #1 & \textbf{\hsmaqasart} & QAS & \textbf{\hsmaqasprio} & #2 \\
				\midrule
				\multicolumn{2}{l}{\textbf{\hsmaqastitel}} & \multicolumn{4}{l}{\parbox[t]{11.8cm}{#3}} \\
				\multicolumn{2}{l}{\textbf{\hsmaqasquelle}} & \multicolumn{4}{l}{\parbox[t]{11.8cm}{#4}} \\
				\multicolumn{2}{l}{\textbf{\hsmaqasstimulus}} & \multicolumn{4}{l}{\parbox[t]{11.8cm}{#5}} \\
				\multicolumn{2}{l}{\textbf{\hsmaqasartefakt}} & \multicolumn{4}{l}{\parbox[t]{11.8cm}{#6}} \\
				\addlinespace
				\multicolumn{6}{l}{\textbf{\hsmaqasumgebung}} \\
				\multicolumn{6}{l}{\parbox[t]{13.5cm}{#7\strut}} \\
				\addlinespace
				\multicolumn{6}{l}{\textbf{\hsmaqasantwort}} \\
				\multicolumn{6}{l}{\parbox[t]{13.5cm}{#8\strut}} \\
				\addlinespace
				\multicolumn{6}{l}{\textbf{\hsmaqasmass}} \\
				\multicolumn{6}{l}{\parbox[t]{13.5cm}{#9\strut}} \\
				\bottomrule
			\end{tabularx}
		\end{footnotesize}
	\end{table}
}


% -------------------------------------------------------
% Planungsnotizen: Sobald ein Abschnitt ausgeschrieben ist, die \plan{...}-Box
% dort löschen. Alle auf einmal ausblenden: \plantrue durch \planfalse ersetzen.
\newif\ifplan
\plantrue
\newcommand{\plan}[1]{%
  \ifplan
    \par\noindent\colorbox{black!6}{\parbox{\dimexpr\linewidth-2\fboxsep\relax}{%
      \small\sffamily\textbf{Geplant:}\ #1}}\par\medskip
  \fi}

% Kennzeichnung, wer einen Abschnitt verfasst hat (direkt nach der Überschrift)
\newcommand{\verfasser}[1]{%
  \par\vspace{-0.6\baselineskip}\noindent{\small\sffamily\color{black!55}Verfasst von: #1}\par\medskip}

% Grafiken
\usetikzlibrary{arrows.meta,positioning,fit,calc,shapes.geometric,backgrounds,decorations.pathreplacing,shapes.multipart}

% Platzhalter für Screenshots, bis echte Bilder vorliegen
\newcommand{\screenshotplatzhalter}[2]{%
  \fbox{\parbox[c][#1][c]{0.9\linewidth}{\centering\sffamily\small\color{black!55}Screenshot folgt:\\ #2}}}

% Screenshot mit dünnem Rahmen (die App hat einen weißen Hintergrund)
\newcommand{\screenshot}[2][\linewidth]{{\setlength{\fboxsep}{0pt}\setlength{\fboxrule}{0.4pt}\fbox{\includegraphics[width=#1]{screenshots/#2}}}}

% Platzhalter für noch fehlende Quellen. Ausgabe: [X]. Das optionale Argument
% beschreibt, welche Quelle gemeint ist, und erscheint nicht im PDF.
% Beispiel: \X[RFC 6455, Abschnitt 4]
\newcommand{\X}[1][]{\mbox{[X]}}

% Syntax für JavaScript/TypeScript-Listings
\lstdefinelanguage{JavaScript}{
  keywords={const,let,var,function,return,if,else,for,of,in,new,import,from,export,async,await,try,catch,throw,switch,case,break,default,typeof,null,true,false,continue},
  sensitive=true,
  morecomment=[l]{//},
  morecomment=[s]{/*}{*/},
  morestring=[b]",
  morestring=[b]',
  morestring=[b]`
}

% Mehr Spielraum beim Umbruch langer Code-Bezeichner
\setlength{\emergencystretch}{3em}

% Markierungen fuer die Ueberarbeitung (vor der Abgabe entfernen)
\usepackage[most]{tcolorbox}
\newtcolorbox{neu}{breakable, enhanced, colback=yellow!40, colframe=yellow!40, boxrule=0pt, arc=0pt, left=3pt, right=3pt, top=2pt, bottom=2pt}
\newtcolorbox{offen}{breakable, enhanced, colback=yellow!40, colframe=orange!85!black, boxrule=1.2pt, arc=0pt, left=3pt, right=3pt, top=2pt, bottom=2pt, title={Noch klären}, fonttitle=\sffamily\small\bfseries, coltitle=black, colbacktitle=orange!45}
\usepackage{soul}
\sethlcolor{yellow!40}
\newcommand{\neutext}[1]{\hl{#1}}
